% Page 064

\seite{64}

Die diesjährigen Herbstferien dauerten vom 22.\ September\ob
bis einschließlich 21.\ Oktober.

Nach einigem schönen, trockenen Wetter herrscht\ob
heuer durchschnittlich Kälte und Nässe vor, so daß das\ob
Getreide auf den Feldern sehr darunter zu leiden hat.\ob
Großen Schaden leiden vor allen Dingen die Kartoffeln,\ob
von denen 40--50 % in Fäulnis übergehen.

Ein Fortschritt der Zeit ist auch in hiesiger Gemeinde\ob
zu verzeichnen. Laut Gemeinderats-Beschluß soll nämlich\ob
die elektrische Lichtanlage errichtet werden.

\margmark{15. Dezember.}%
Die Weihnachtsferien dauern vom 17.\ Dezember\ob
bis einschließlich 7.\ Januar.

\pend
\abschnitt{1923}
\pstart

Jahre kommen, Jahre vergehen. Das neue hat be\ohb
gonnen, das Alte ist zerronnen. Ein Blick auf die\ob
augenblickliche Welt- \& Marktlage sagt uns, wo wir\ob
stehen – vor einem tiefen Abgrund, über den 1922 keine\ob
Brücken zu bauen vermochte, sondern uns seinem Rande\ob
nur einige Schritte näher brachte. Die ersehnte und er\ohb
hoffte Hilfe blieb aus. Wie wird es weiter gehen? Unsere\ob
Hoffnung schwindet. Vom Westen ziehen neue unheil\ohb
volle Gewitterwolken empor. Blitze zucken schon auf das\ob
Ruhrgebiet und die friedlichen Rheinlande. Wo werden\ob
sie ihr Opfer verschlingen?! Laßt uns die Hoffnung fest\ohb
halten, Gott wird uns wunderbar erhalten!!!

\margmark{23. Januar.}%
Heute vormittag fand eine Revision des Gesundheits\ohb
zustandes der Schüler durch den Herren Kreisarzt aus\ob
Bitburg statt.

Die politische Lage wird immer verwickelter.\ob
Frankreich und Belgien haben das Ruhrgebiet und Teile\ob
darüber hinaus besetzt. Demzufolge hat die\ob
deutsche Regierung die Reparationslieferungen eingestellt.\ob
Die Franzosen erteilen Befehle, die deutsche Regierung\ob
erteilt Gegenbefehle. Was wird daraus werden?

\pend
\abschnitt{fr. 13.I.23. W}
\pstart

\pend
